%! Linear Functions: Slope, Forms & Graphing — Unit 1, Lesson 1.1 — Guided Notes
% Gradual release (user direction, 2026-08-31): the notes are the "I do / we do" block, 20 minutes,
% and the Individual Practice block that closes them is the "you do alone", 15 more.
\documentclass[10pt]{article}
\usepackage{algebra2-article}
\usepackage{algebra2-boxes}
\newcommand{\TallMath}[1]{$\displaystyle #1\rule[-1.4em]{0pt}{3.2em}$}
% Coordinate grid: {xmin}{xmax}{ymin}{ymax}
\newcommand{\lgrid}[4]{%
  \draw[step=1, linegray!40, very thin] (#1,#3) grid (#2,#4);
  \draw[->, semithick, charcoal] (#1,0)--(#2,0) node[below right, font=\scriptsize]{$x$};
  \draw[->, semithick, charcoal] (0,#3)--(0,#4) node[left, font=\scriptsize]{$y$};}

\begin{document}

\pageheader{Unit 1, Lesson 1.1}{Guided Notes}

\vspace{0.10in}

\begin{objectivebox}
\textbf{By the end of this lesson, I will be able to\dots}
\begin{itemize}\setlength{\itemsep}{2pt}
  \item find the \emph{slope} of a line from a graph or from \emph{any} two points;
  \item write a line in \emph{slope-intercept}, \emph{point-slope}, and \emph{standard} form, and
        convert among them --- even when I am never told the starting value;
  \item describe a line as a transformation of the \emph{parent} $y=x$;
  \item identify \emph{parallel}, \emph{perpendicular}, horizontal, and vertical lines.
\end{itemize}
\end{objectivebox}

\vspace{0.05in}

\begin{vocabbox}
\small
Fill in each term as we build it together.
\par\vspace{2pt}   % \par is required: \termblanklong's \noindent is a no-op mid-paragraph
\termblanklong{Slope from two points}
\termblanklong{Slope-intercept form}
\termblanklong{Point-slope form}
\termblanklong{Standard form}
\termblanklong{Parent function $y=x$ / transformation}
\termblanklong{Parallel / perpendicular}
\end{vocabbox}

\begin{hookbox}
Your phone bill does not list the monthly fee --- only the total. Two months of bills:
\textbf{4~GB cost \$52}, and \textbf{10~GB cost \$70}.
\begin{itemize}\setlength{\itemsep}{1pt}
  \item From one bill to the other: \blank{1.2cm} more GB for \$\blank{1.2cm} more, so each
        GB costs \$\blank{1.0cm}.
  \item So what is the flat monthly fee --- the cost of \emph{zero} GB? \blank{1.4cm}
\end{itemize}
Nobody handed you the starting value; you \emph{worked back} to it. Today we build the equation of a
line from whatever we are given --- a graph, two points, or a slope and \emph{any} point at all.
\end{hookbox}

\boxguard[22]
\begin{notesbox}{1. Slope from Any Two Points}
\begin{center}
\begin{minipage}[c]{0.50\linewidth}
\[ m \;=\; \frac{\Delta y}{\Delta x} \;=\; \frac{y_2-y_1}{x_2-x_1} \]
Subtract in the \emph{same order} top and bottom. \emph{Any} two points on the line give the same
answer, because a line's rate of change is \blank{2.2cm}.
\end{minipage}\hfill
\begin{minipage}[c]{0.46\linewidth}\centering
\begin{tikzpicture}[scale=0.42]
  \lgrid{-3}{5}{-6}{6}
  \draw[very thick, forest] (-2,-6)--(4,6);
  \draw[semithick, navy, dashed] (1,0)--(2,0)--(2,2);
  \node[navy, font=\tiny] at (1.5,-0.8){run};
  \node[navy, font=\tiny, right] at (2.1,1){rise};
  \fill[charcoal] (0,-2) circle (3.4pt);
  \fill[charcoal] (3,4) circle (3.4pt);
\end{tikzpicture}
\end{minipage}
\end{center}
\textbf{Worked: the phone plan.} The bills give the points $(4,52)$ and $(10,70)$.
\begin{work}
  m &= \frac{70 - 52}{10 - 4} \\
    &= \frac{18}{6} \\
    &= 3
\end{work}
So the plan costs \$\blank{1.0cm} per GB.
\end{notesbox}

\boxguard[24]
\begin{notesbox}{2. Three Ways to Write One Line}
\begin{center}
\renewcommand{\arraystretch}{1.5}
\begin{tabularx}{\linewidth}{@{}l l X@{}}
\textbf{Form} & \textbf{Looks like} & \textbf{Hands you \dots\ / use it when \dots} \\
\hline
Slope-intercept & $y = mx + b$ & the \blank{2.0cm} and the \blank{2.4cm} \\
Point-slope & $y - y_1 = m(x - x_1)$ & a slope and \blank{3.0cm} \\
Standard & $Ax + By = C$ & integers, $A \ge 0$; tidy for reading \blank{2.0cm} \\
\end{tabularx}
\end{center}
\textbf{Why point-slope exists.} In the phone plan you had a rate and a point --- but the
$y$-intercept was the one number nobody gave you. Point-slope lets you write the line
\emph{immediately}, from the point you already have:
\[ y - \blank{1.0cm} \;=\; \blank{1.0cm}\,(x - \blank{1.0cm}) \qquad
   \text{using the point } (4,52). \]
\end{notesbox}

\boxguard[24]
\begin{notesbox}{3. Same Line, Different Outfit}
Two equations that \emph{look} different can be the same line. Expand point-slope to get
slope-intercept:
\begin{work}
  y - 52 &= 3(x - 4) \\
  y - 52 &= 3x - 12 \\
       y &= 3x + 40
\end{work}
There is the monthly fee: $b = \blank{1.2cm}$ --- the same number you worked back to in the hook.
Now move $x$ across to get standard form:
\begin{work}
  y &= 3x + 40 \\
  -3x + y &= 40 \\
  3x - y &= -40
\end{work}
\textbf{Test any conversion} by putting one input through both rules. At $x=10$: \; slope-intercept
gives \blank{1.4cm}, and the original bill said \blank{1.4cm}.
\end{notesbox}

\boxguard[24]
\begin{notesbox}{4. The Parent $y=x$, and Four Special Cases}
Every non-vertical line is the parent $y=x$ after \emph{two} moves.
\begin{center}
\begin{minipage}[c]{0.46\linewidth}
\begin{itemize}[leftmargin=1.4em, itemsep=2pt]
  \item Multiplying by $k$ changes the \blank{1.8cm} --- steeper, flatter, or (if $k<0$)
        \blank{1.6cm}.
  \item Adding $k$ changes the \blank{2.4cm}: the line slides up or down.
\end{itemize}
\end{minipage}\hfill
\begin{minipage}[c]{0.50\linewidth}\centering
\begin{tikzpicture}[scale=0.40]
  \lgrid{-4}{5}{-6}{6}
  \draw[very thick, navy, dashed] (-4,-4)--(5,5) node[above left, font=\tiny, navy]{$y=x$};
  \draw[very thick, forest] (-1,-6)--(5,6) node[below right, font=\tiny, forest]{$y=2x-4$};
  \fill[charcoal] (0,-4) circle (3.4pt);
\end{tikzpicture}
\end{minipage}
\end{center}
\begin{center}
\renewcommand{\arraystretch}{1.4}
\begin{tabularx}{\linewidth}{@{}l X l X@{}}
\textbf{Parallel} & slopes are \blank{1.8cm} & \textbf{Perpendicular} & slopes are \blank{3.0cm}, so $m_1m_2 = \blank{1.0cm}$ \\
\textbf{Horizontal} $y=c$ & slope \blank{1.0cm} & \textbf{Vertical} $x=a$ & slope \blank{2.0cm} \\
\end{tabularx}
\end{center}
Only one of those four \emph{cannot} be written as $y = mx+b$: \blank{2.0cm} --- because it has
no slope to put in for $m$.
\end{notesbox}

\begin{practicebox}
A line has slope $-\tfrac12$ and passes through $(2,-5)$.
\begin{enumerate}[leftmargin=1.7em, itemsep=3pt]
  \item \textbf{Point-slope} form: \blank{4.6cm}
  \item Convert to \textbf{slope-intercept} form.
        \begin{work}
          y + 5 &= -\tfrac12(x - 2) \\
          y + 5 &= -\tfrac12 x + 1 \\
              y &= -\tfrac12 x - 4
        \end{work}
  \item A line \textbf{parallel} to it has slope \blank{1.4cm}; a line \textbf{perpendicular}
        to it has slope \blank{1.4cm}.
\end{enumerate}
\end{practicebox}

\boxguard[22]
\begin{scenariobox}[Individual Practice --- On Your Own]{navy}
Three problems, quietly and on your own. Every one of them is a line you have to \emph{write} ---
from two points, from a form you were handed, or from a rate and a start.
\begin{enumerate}[leftmargin=1.9em, itemsep=9pt]
  \item \textbf{From two points.} A line passes through $(-1,4)$ and $(3,-4)$.
        \begin{enumerate}[label=(\alph*), leftmargin=1.9em, itemsep=3pt]
          \item Slope: \blank{1.4cm}
          \item Point-slope form, using $(3,-4)$: \blank{4.4cm}
          \item Now convert to slope-intercept form.
            \begin{work}
              y + 4 &= -2(x - 3) \\
              y + 4 &= -2x + 6 \\
                  y &= -2x + 2
            \end{work}
        \end{enumerate}

  \item \textbf{From standard form.} Rewrite $3x - 4y = 12$ in slope-intercept form.
    \begin{work}
      3x - 4y &= 12 \\
         -4y &= -3x + 12 \\
            y &= \tfrac34 x - 3
    \end{work}
    Slope: \blank{1.4cm} \quad $y$-intercept: \blank{1.6cm}

  \item \textbf{From a rate and a start.} A line has slope $\tfrac23$ and crosses the $y$-axis at
        $(0,5)$. Write it in slope-intercept form, then in \textbf{standard} form
        (integers, $A \ge 0$).
    \begin{work}
          y &= \tfrac23 x + 5 \\
         3y &= 2x + 15 \\
    -2x + 3y &= 15 \\
     2x - 3y &= -15
    \end{work}
    Slope-intercept: \blank{3.2cm} \quad Standard: \blank{3.0cm}
\end{enumerate}
\end{scenariobox}

\end{document}
